% ============================================================
%  1. Introduction
% ============================================================
\section{Introduction}
\label{sec:intro}

本文以课程大作业中的A股日频选股任务为背景，并采用research-style technical report的组织方式。为了让报告保持可读性，本文始终围绕一个核心问题展开：\textbf{能否学习一个稳定的横截面选股信号，使模型每天把更值得买入的股票排到前面？} 在全市场约5000只股票中每日选出top-20持仓，本质上是一个横截面排序问题：投资决策主要取决于股票之间的相对顺序，而非对单只股票次日收益的精确回归。

需要强调的是，本文研究的是\textbf{信号层}而非完整\textbf{交易执行层}。我们的离线评估协议使用top-20等权日调仓来横向比较不同模型的排序能力；该协议有利于检验模型是否能在holdout集上稳定选出相对更优的股票，但并未完整建模手续费、滑点、涨跌停成交约束、换手率控制和多日持有优化。因此，本文报告的离线收益指标应理解为受控协议下的选股信号表现，而不是可直接落地的真实净收益策略。为了检验从模型信号到真实操作之间的差距，本文还记录了一次同花顺模拟交易；该模拟结果较差，但它衡量的是“日频alpha信号+人工盘中执行”的端到端过程，而不是CLIME纯排序信号本身。

我们的核心诊断有两点。第一，股票预测的训练目标需要更贴近最终的排序决策。一方面，学习排序和偏好学习中的经验表明，pairwise supervision适合学习相对判断~\cite{joachims2002optimizing,burges2005learning,rendle2009bpr,christiano2017deep,ouyang2022training}；另一方面，单纯回归损失会对所有数值误差施加近似对称惩罚，却无法区分“方向错误”和“幅度略有偏差”这两类在选股中影响完全不同的错误。第二，只看单只股票自己的历史序列通常不够：A股中的行业联动、风格轮动和市场状态会同时影响一批股票。若模型完全看不到同业和市场上下文，其上限会受到明显限制。

因此，本文最重要的设计不是简单堆叠更大的Transformer，而是引入\textbf{peer-aware market context}：对每只股票构造top-10相似同业的动态统计量，并与市场上下文一起注入模型。市场信息的注入也需要谨慎处理：简单拼接市场指数或同业信息可能把噪声直接送入Backbone，而不是形成对个股判断的有约束修正。CLIME中的ScaledGatedEncoder正是为了解决这一点：它只学习一个受门控和幅度约束的表征偏移，让peer dynamics作为增量修正进入模型。

基于这些观察，我们提出\textbf{CLIME}（Cross-Modal Injection via Learned Market Encoding），一种面向股票排序的Transformer架构。CLIME包含三项关键设计：(1)~Stage~1使用pairwise ranking loss预训练Backbone，先学习通用排序先验；(2)~Stage~2通过BCE方向预热和Directional Regression Loss逐步过渡到次日收益预测；(3)~ScaledGatedEncoder将同业动态和市场上下文以受控加法方式注入表征空间，使市场信息作为增量修正而非直接覆盖个股判断。主实验均使用纯alpha分数（$\lambda=0$）评估CLIME本身的排序能力；风险过滤器仅作为模拟交易中的工程补充。

这种拆分也决定了本文对模拟交易结果的解释方式。若离线排序模型表现较好，但盘中模拟交易收益较差，不能简单得出“模型没有学到信号”的结论。真实收益还受到买入时点、卖出顺序、资金释放、仓位集中、人工择时和市场日内路径的影响。换言之，最终交易收益可以被视为alpha信号、组合构造、执行价格和人工干预的叠加结果。第~\ref{sec:trading}节将展示一个失败的模拟交易案例，并将其作为执行层缺口分析，而不是作为CLIME离线排序能力的主要评价指标。

在holdout测试集（2025年12月至2026年5月，共100个交易日）上，CLIME取得\textbf{累计收益+137.15\%、超额收益+123.14\%}（Sharpe 6.94），相对最强基线Ridge（+91.65\%超额）高出约31个百分点。消融实验显示，Stage~1预训练、BCE课程预热、Directional Regression Loss和ScaledGatedEncoder的移除都会造成明显退化；内部分析进一步表明，模型主要依赖横截面排名、基本面等信号，并形成了紧凑的主导排序方向。

本文的贡献可以概括为四点：第一，将A股日频选股明确建模为top-K排序任务，并清楚区分选股信号评估与真实交易执行；第二，设计了peer dynamics驱动的受控市场信息注入机制，而不是只依赖单股历史序列；第三，通过基线对比、消融实验、Jacobian/SVD、线性探针、特征重要性和极端行情切片，同时评估模型收益、组件贡献和可解释性；第四，结合模拟交易记录分析离线alpha信号到真实执行之间的断层，说明一个表现较好的predictor仍需要独立的portfolio construction和execution模块才能成为完整量化系统。
